\section{Frozen mathematical scope}
The game is linear maximal-run deletion. Circular chord diagrams do not
change the move semantics. The counting objects below are certificate
partitions of a fixed word, not winning words or deletion histories.
Proof snapshot: \texttt{3446f63}, Lean 4.19.0.
The manuscript and proof ledger are available from the project home page.
\section{Definitions and checked results}
\begin{definition}[Color pattern]\label{def:pattern}
A color pattern on $n$ positions is an equivalence relation $p$ on
$[n]=\{0,\ldots,n-1\}$. For a coloring $w:[n]\to\Sigma$, put
\[
 i\mathrel p j\quad\Longleftrightarrow\quad w_i=w_j.
\]
Write $[i]_p=\{j:i\mathrel p j\}$. A pattern is winning when one, equivalently
every, realization of its color equalities is winning.
\lean{LC004.ColorPattern}
\leanok
\end{definition}

\begin{definition}[Noncrossing]\label{def:noncrossing}
A partition $c$ of $[n]$ is noncrossing if, whenever
$a<b<c'<d$ with $a\sim_c c'$ and $b\sim_c d$, one also has $a\sim_c b$.
Thus interleaving pairs may occur only inside a single block.
\lean{LC004.PatternNoncrossing}
\leanok
\end{definition}

\begin{definition}[Certificate partition]\label{def:certificate}
For a pattern $p$, let
\[
 \Cert(p)=\{c:\ c\text{ is a noncrossing partition of }[n],\
              \card{B}\in\{2,3\}\text{ for every block }B,\
              c\leq p\},
\]
where $c\leq p$ means $a\sim_c b\Rightarrow a\mathrel p b$.
\lean{LC004.CertificatePartition}
\leanok
\end{definition}

The imported certificate characterization is
\begin{equation}\label{eq:certificate_characterization}
 \Win(p)\quad\Longleftrightarrow\quad\Cert(p)\ne\varnothing.
\end{equation}

\begin{definition}[Rigid and heavy]\label{def:rigid_heavy}
A position $i$ is rigid if $\card{[i]_p}\in\{2,3\}$ and heavy if
$\card{[i]_p}\geq4$. Let $H(p)$ be the set of heavy positions.
A singleton position is neither rigid nor heavy.
\lean{LC004.RigidPosition,LC004.HeavyPosition}
\leanok
\end{definition}

\begin{lemma}[Local rigidity]\label{lem:rigidity}
Let $c\in\Cert(p)$ and let $i$ be rigid. The $c$-block containing $i$ is
the entire class $[i]_p$.
\lean{LC004.certificate_rigid_block_eq_class}
\leanok
\uses{def:certificate,def:rigid_heavy}
\end{lemma}
\begin{proof}
The block is monochromatic, so it is contained in $[i]_p$, and it has
cardinality two or three. A proper inclusion can only leave one position
uncovered in a class of size three. That remaining position would have to
belong to a different monochromatic certificate block containing at least two
positions. Exact coverage makes the two certificate blocks disjoint, so the
second block would have to fit into a singleton. This is impossible.
\leanok
\uses{def:certificate}
\end{proof}

\begin{definition}[Face signature]\label{def:face}
Put $a\sim_F b$ when
\begin{equation}\label{eq:face_signature}
 \forall u,v,\quad
 u\text{ rigid and }u\mathrel p v
 \ \Longrightarrow\
 (u<a<v\Longleftrightarrow u<b<v).
\end{equation}
For an anchor $a$, its complete heavy carrier is
\[
 H_a=\{x\in[n]:x\in H(p),\ a\sim_F x\}.
\]
The increasing enumeration is an order embedding
$e_a:[\card{H_a}]\hookrightarrow[n]$.
\lean{LC004.rigidFacePattern,LC004.heavyFacePositions,LC004.heavyFaceEmbedding}
\leanok
\uses{def:rigid_heavy}
\end{definition}

\begin{lemma}[Arc criterion]\label{lem:arc}
For heavy positions $a<b$, one has $a\sim_F b$ if and only if there are no
same-colored rigid positions $u,v$ such that
\[
 a<u<b,\qquad v<a\ \text{or}\ b<v.
\]
\lean{LC004.rigid_face_iff_no_rigid_arc_separator}
\leanok
\uses{def:face}
\end{lemma}
\begin{proof}\leanok
The complete proof is in the checked Lean declarations linked above.
\end{proof}


\begin{theorem}[Heavy-block confinement]\label{thm:confinement}
If $c\in\Cert(p)$, $a$ is heavy, and $a\sim_c b$, then $b$ is heavy and
$a\sim_F b$.
\lean{LC004.certificate_heavy_face}
\leanok
\uses{lem:rigidity,def:face}
\end{theorem}
\begin{proof}
Monochromaticity gives $a\mathrel p b$, so $a$ and $b$ have the same original
color-class size. Thus $b$ is heavy. Fix a rigid chord $u,v$ of one color.
By Lemma~\ref{lem:rigidity}, $u$ and $v$ lie in one forced certificate block.
If precisely one of $a,b$ lies in $(u,v)$, the two certificate blocks have
alternating endpoint pairs in one of their orders. Their four endpoints are
distinct because heavy and rigid positions are disjoint. This contradicts
noncrossing for blocks of size at most three. Therefore the interval tests
agree for every rigid chord.
\leanok
\uses{lem:rigidity,def:noncrossing}
\end{proof}

\begin{corollary}[Face saturation]\label{lem:saturation}
If $e$ enumerates all and only the heavy positions in one face and
$c(e(i),b)$, then $b$ is in the range of $e$.
\lean{LC004.certificate_face_saturated}
\leanok
\uses{thm:confinement}
\end{corollary}
\begin{proof}\leanok
The complete proof is in the checked Lean declarations linked above.
\end{proof}


Write $E_{c,d}(x,y)$ for $c(x,y)\lor d(x,y)$. Restriction is the pullback
$\Res_e(c)(i,j)\Longleftrightarrow c(e(i),e(j))$; $c\vee d$ is the least
equivalence relation containing both relations.
\begin{definition}[Path locality]\label{def:locality}
A carrier $e$ has path locality for $c,d$ if every global $E_{c,d}$-path between
$e(i)$ and $e(j)$ can be represented by an internal path in the carrier
between $i$ and $j$.
\lean{LC004.FacePathLocality}
\leanok
\end{definition}

\begin{lemma}[Locality from saturation]\label{lem:locality}
If $e$ is injective and its range is closed under every $c$-edge and every
$d$-edge, then $e$ has path locality for $c,d$.
\lean{LC004.facePathLocality_of_edge_saturation}
\leanok
\uses{def:locality}
\end{lemma}
\begin{proof}
Induct on the number of edges in a global path. Lift its initial vertex using
the given starting index. If the prefix ends at $e(k)$ and the next edge
reaches $z$, saturation supplies an index $\ell$ with $e(\ell)=z$.
Append the corresponding restricted edge. At the final vertex,
injectivity identifies the constructed index with the prescribed endpoint.
The zero-edge path gives reflexivity.
\leanok
\end{proof}

\begin{theorem}[Geometric join restriction]\label{thm:join_restriction}
For two certificates $c,d$ of the same color pattern and every canonical
complete heavy face,
\[
 \Res_{e_a}(c\vee d)=\Res_{e_a}(c)\vee\Res_{e_a}(d).
\]
\lean{LC004.canonical_heavy_face_restrict_join}
\leanok
\uses{lem:saturation,lem:locality}
\end{theorem}
\begin{proof}\leanok
The complete proof is in the checked Lean declarations linked above.
\end{proof}


\begin{lemma}[Restriction preserves winning]\label{lem:restriction}
If $p$ is winning, then every restricted pattern $\Res_{e_a}(p)$ is winning.
The increasing word in the original colors on $H_a$ is erasable by maximal
runs.
\lean{LC004.winning_restrict_heavy_face,LC004.raw_win_heavy_face_word}
\leanok
\uses{lem:saturation}
\end{lemma}
\begin{proof}
Choose a global certificate. Saturation says that every certificate class
meeting the face is wholly contained in it. Consequently restriction
preserves its cardinality, so the restricted classes still have size two
or three. The increasing embedding preserves the ordering of four endpoints
and hence noncrossing. Monochromaticity is preserved by pullback.
Equation~\eqref{eq:certificate_characterization} gives the local win, and the
imported semantic correspondence returns a maximal-run erasure in the
original face colors.
\leanok
\end{proof}

\begin{theorem}[Structural Face Lemma]\label{thm:face_lemma}
For every finite color pattern $p$,
\begin{align*}
 \Win(p)\quad\Longleftrightarrow\quad&
 \bigl(\forall i,\ \card{[i]_p}\geq2\bigr)\\
 &{}\land\bigl(\text{the rigid skeleton has no alternating chord pairs}\bigr)\\
 &{}\land\bigl(\forall a,\ \Win(\Res_{e_a}(p))\bigr).
\end{align*}
The same equivalence holds for maximal-run erasure of a word with its
original colors.
\lean{LC004.winning_iff_heavy_face_decomposition,LC004.raw_win_iff_heavy_face_decomposition}
\leanok
\uses{lem:rigidity,lem:restriction,thm:cellular_assembly}
\end{theorem}
\begin{proof}\leanok
The complete proof is in the checked Lean declarations linked above.
\end{proof}


Let $(S_k)_{k\in I}$ be disjoint regions covering $[n]$, with increasing
enumerations $e_k:[m_k]\hookrightarrow[n]$, and suppose
\begin{equation}\label{eq:region_noncrossing}
 e_k(i)<e_\ell(s)<e_k(j)<e_\ell(t)\quad\Longrightarrow\quad k=\ell.
\end{equation}
For certificate partitions $c_k\in\Cert(\Res_{e_k}(p))$, define
\begin{equation}\label{eq:assembled}
 x\sim_g y\quad\Longleftrightarrow\quad
 \exists k,i,j,\ e_k(i)=x,\ e_k(j)=y,\ i\sim_{c_k}j.
\end{equation}
\begin{lemma}[Assembly on regions]\label{lem:region_assembly}
The relation $g$ is a certificate partition of $p$, and
$\Res_{e_k}(g)=c_k$ for every $k$.
\lean{LC004.assembledPattern_constraint,LC004.exists_partition_of_noncrossing_regions}
\leanok
\uses{def:certificate}
\end{lemma}
\begin{proof}
Coverage supplies a region and local coordinate for each position, proving
reflexivity. Symmetry is inherited locally. For transitivity, two witnesses
sharing a middle position must belong to the same region by disjointness.
Injectivity of that region's embedding identifies their middle coordinates,
and local transitivity completes the argument.

A $g$-class meeting $S_k$ stays inside $S_k$. Its pullback is precisely the
corresponding $c_k$-class, so its cardinality is two or three. If two
$g$-relations have alternating endpoints, \eqref{eq:region_noncrossing} puts
them in the same region. The increasing embedding reflects their inequalities;
local noncrossing then implies the required relation between the first two
endpoints. Finally, the local refinement $c_k\leq\Res_{e_k}(p)$ implies
$g\leq p$.

To verify exact restriction, unfold \eqref{eq:assembled}. A witness connecting
two points of $S_k$ cannot come from a different region. Injectivity then
recovers the original local coordinates. The converse uses $k$ itself.
\leanok
\end{proof}

Let $Q_F=[n]/{\sim_F}$ and choose the quotient representative used by Lean.
Write $e_q$ for the increasing embedding of its complete heavy carrier.
The restriction map is
\[
 R:\Cert(p)\longrightarrow\prod_{q\in Q_F}\Cert(\Res_{e_q}(p)),
 \qquad R(c)_q=\Res_{e_q}(c).
\]
\begin{theorem}[Cellular assembly isomorphism]\label{thm:cellular_assembly}
Suppose every appearing color has at least two positions and the rigid
skeleton is noncrossing. Then $R$ is a bijection:
\[
 \Cert(p)\ \simeq\ \prod_{q\in Q_F}\Cert(\Res_{e_q}(p)).
\]
\lean{LC004.heavyFaceCertificateEquiv}
\leanok
\uses{thm:confinement,lem:rigidity,lem:region_assembly}
\end{theorem}
\begin{proof}
For injectivity, suppose two global partitions have the same face
restrictions. A relation beginning at a rigid position is forced by the
original color class, so it holds in both partitions. A relation beginning
at a heavy position stays in one heavy face by confinement. Lift its endpoints
into the increasing face coordinates, use equality of restrictions, and
return to the global positions. Interchanging the two partitions gives
equality.

For surjectivity, take any tuple of local partitions. Supply the forced
rigid blocks and apply the regional assembly. Exact restriction in
Lemma~\ref{lem:region_assembly} proves that the resulting global partition
maps to the supplied tuple. Thus assembly recovers every tuple, and
injectivity shows that it does so uniquely.
\leanok
\end{proof}

\begin{corollary}[Certificate product]\label{thm:certificate_product}
Under the hypotheses of Theorem~\ref{thm:cellular_assembly},
\[
 \card{\Cert(p)}
 =\prod_{q\in Q_F}\card{\Cert(\Res_{e_q}(p))}.
\]
\lean{LC004.certificate_partition_count_product}
\leanok
\uses{thm:cellular_assembly}
\end{corollary}
\begin{proof}
Every partition type in the statement is finite. Apply invariance of finite
cardinality under equivalence and the cardinality formula for a finite
dependent product.
\leanok
\end{proof}\section{Open publication obligations}
\begin{proposition}[Ordered slots]\label{prop:slots}\notready
For fixed face words and ordered gap carriers, prove the weak-composition
placement bijection in Lean, including the extra outer slot.
\end{proposition}
The complete ordinary proof is in the manuscript. Counterexample minimality
is computationally established; no Lean declaration is attached to it.
Transaction and streaming contracts are architecture specifications and
are not kernel-verified theorems about Alexandria Core.
